\section{Historical SIR:
BIRA}\label{record-docs-watermarking-qualitative-paper-sir-bira-md-historical-sir-bira}

Historical OPT/Llama control. Fixed positions are drawn from the first
330 responses available in all arms. Missing owner or quality scores
remain explicitly unmeasured.

Eight pairs selected by fixed test positions before inspecting detection
or quality outcomes. Undetected originals and unsuccessful attacks
remain included.

\begin{recordfigure}
\centering
\pandocbounded{\includegraphics[keepaspectratio,alt={Sixteen-panel text preview}]{figures/paper-sir-bira.png}}
\caption{Sixteen-panel text preview}
\end{recordfigure}

The figure previews up to the first 360 original characters per
response; full texts below are unabridged. Emoji use Unicode codepoints
in the preview to avoid missing-glyph boxes. Owner scores and quality
metrics use full responses.

{\scriptsize\setlength{\tabcolsep}{3pt}\def\LTcaptype{none} % do not increment counter
\begin{longtable}[]{@{}
  >{\raggedright\arraybackslash}p{(\linewidth - 8\tabcolsep) * \real{0.2000}}
  >{\raggedright\arraybackslash}p{(\linewidth - 8\tabcolsep) * \real{0.2000}}
  >{\raggedright\arraybackslash}p{(\linewidth - 8\tabcolsep) * \real{0.2000}}
  >{\raggedright\arraybackslash}p{(\linewidth - 8\tabcolsep) * \real{0.2000}}
  >{\raggedright\arraybackslash}p{(\linewidth - 8\tabcolsep) * \real{0.2000}}@{}}
\toprule\noalign{}
\begin{minipage}[b]{\linewidth}\raggedright
Example
\end{minipage} & \begin{minipage}[b]{\linewidth}\raggedright
Test position
\end{minipage} & \begin{minipage}[b]{\linewidth}\raggedright
Before
\end{minipage} & \begin{minipage}[b]{\linewidth}\raggedright
After
\end{minipage} & \begin{minipage}[b]{\linewidth}\raggedright
Jina
\end{minipage} \\
\midrule\noalign{}
\endhead
\bottomrule\noalign{}
\endlastfoot
1 & 0 & Owner score not measured & Owner score not measured & not
measured \\
2 & 41 & Owner score not measured & Owner score not measured & not
measured \\
3 & 82 & Owner score not measured & Owner score not measured & not
measured \\
4 & 123 & Owner score not measured & Owner score not measured & not
measured \\
5 & 165 & Owner score not measured & Owner score not measured & not
measured \\
6 & 206 & Owner score not measured & Owner score not measured & not
measured \\
7 & 247 & Owner score not measured & Owner score not measured & not
measured \\
8 & 288 & Owner score not measured & Owner score not measured & not
measured \\
\end{longtable}
}

\subsection{Example 1: test position
0}\label{record-docs-watermarking-qualitative-paper-sir-bira-md-example-1-test-position-0}

Prompt ID:
\texttt{\seqsplit{e31e6f74ee79d4c901f4c7e2d71bc20b7d4ac16a960885a84dbed35ed994c04a}}.

\textbf{Question / prompt}

\begin{Verbatim}
Cluster comprises IBM's Opteron-based eServer 325 server and systems management software and storage devices that can run Linux and Windows operating systems
\end{Verbatim}

\textbf{Original watermarked-model output}

\begin{Verbatim}
Cluster comprises IBM's Opteron-based eServer 325 server and systems management software and storage devices that can run Linux and Windows operating systems. The system manages virtualized and real virtualized hardware with two server clusters on different platforms, each with 1,048 Dell Hyperserver server nodes, and manages server space and CPU load for all three clusters. Deployed in virtualized and remote distributed datacenters, and in the cloud, Clusters provide cloud and data center scalability along with improved availability and fault tolerance to allow for large-scale deployment of applications and virtual computing assets.

The eServer 325 system was developed to increase efficiency and scalability to satisfy the needs of cloud and datacenter developers.

"Our customers have been requesting virtualized and non-virtualized solutions for cloud and multi-node deployments from IBM for many products and we have seen tremendous market interest," said Bryan Murphy, senior program manager, Clusters at IBM Global Software. "Customers are increasingly looking at Clusters to provide increased capability, control and flexibility to satisfy their applications and data sets."

Availability and pricing of IBM Clusters are based on availability with an enterprise or shared server purchase, including 24-hour virtual machine capacity, as well as shared management, hardware, and
\end{Verbatim}

\textbf{After / alternative output}

\begin{Verbatim}
The cluster is comprised of IBM's eServer 325, a high-performance server, combined with advanced systems management software and scalable storage solutions capable of supporting both Linux and Windows operating systems. This powerful system efficiently handles both physical and virtualized infrastructure, integrating seamlessly with up to two distinct server clusters, each featuring a massive array of 1,048 nodes powered by Dell's Hyperserver technology. Furthermore, it ensures optimal allocation of resources across these multiple clusters by dynamically managing both server spaces and CPU loads. Designed for flexible implementation, these highly adaptable cluster solutions can be seamlessly integrated into both local, remotely hosted, and cloud-based data centers. As such, they offer unparalleled scalability, enhanced reliability, and robust resilience, making it possible to deploy complex, high-capacity applications across vast, geographically dispersed environments. 

The eServer 325 system was specifically designed by IBM engineers to meet the evolving demands of modern cloud and datacenter professionals, focusing primarily on enhanced operational efficiency and streamlined scalability. 

"We have observed a significant surge in demand among our clients seeking innovative, hybrid solutions capable of supporting both virtualized and traditional infrastructure within a single, cohesive framework," stated Bryan Murphy, a seasoned expert within IBM's global software division. "Our clients are actively seeking out IBM's Cluster solutions, driven by a growing need for elevated levels of flexibility, manageability, and adaptability, which are essential for effectively managing their diverse application portfolios and complex datasets."

IBM's Cluster offerings are competitively priced, providing customers access to a comprehensive suite of features, inclusive of dedicated, round-the-clock support, comprehensive resource sharing, as well as collaborative system administration, upon acquiring a qualifying, high-end, commercially licensed server solution.
\end{Verbatim}

\subsection{Example 2: test position
41}\label{record-docs-watermarking-qualitative-paper-sir-bira-md-example-2-test-position-41}

Prompt ID:
\texttt{\seqsplit{f349a8727c7c51a1c91fedf1fcafe0ae6851c5474bb2ce79fccc84c7544358e8}}.

\textbf{Question / prompt}

\begin{Verbatim}
Director Patty Jenkins has revealed details about the upcoming sequel.
If we've learned anything from our era of superhero movies, it's that no one
\end{Verbatim}

\textbf{Original watermarked-model output}

\begin{Verbatim}
Director Patty Jenkins has revealed details about the upcoming sequel.
If we've learned anything from our era of superhero movies, it's that no one makes anything for less money than Marvel Studios does. Having won two Oscars for her brilliant original Wonder Woman, Patty Jenkins made her debut with $750 million and no pressure from Hollywood for two more movies. The good news for a woman who's never seen any comic book characters and not entirely sure if she wants to even, Wonder Woman 1984 wasn't her first foray into comic books, having previously worked with Marc Webb on the highly successful Daredevil and Jessica Jones, with both not being good movies, or movies with any female leads to be frank, and it turned into $790,070,000 at the box office and went on to win an Oscar in its own right for its visual story telling.
When asked if she knows anything about her future in the comic book world, Patty Jenkins has said that not really and if she does want to come back, that she'll certainly be doing her damndest to make her vision known. She has also revealed an interesting bit of information though, saying that she plans her movies with female characters at her heart, and won't be making something for men. Checkout an excerpt from her comments below,
\end{Verbatim}

\textbf{After / alternative output}

\begin{Verbatim}
Patty Jenkins, director of the upcoming sequel, has shared insights about the project. Given the current trend of superhero films, one key aspect stands out: Marvel Studios consistently produces films without exceeding budget expectations. Following the success of her Oscar-winning film 'Wonder Woman,' Patty Jenkins entered the industry with a significant budget of $750 million. Notably, she had the luxury of being under minimal external pressure to produce additional projects. 

For someone who is unfamiliar with the world of comics and hasn't expressed strong enthusiasm about exploring the genre further, the commercial success of 'Wonder Woman 1984' came as a surprise. This film was not Jenkins' first experience working on a comic book adaptation; she had collaborated with Marc Webb on 'Daredevil' and 'Jessica Jones.' Although these projects received mixed reviews and lacked prominent female leads, they still managed to gross $790.07 million at the box office. Furthermore, 'Wonder Woman 1984' received critical acclaim for its storytelling.

During discussions about her involvement in the comic book world, Patty Jenkins expressed uncertainty about her future prospects. However, she is resolute about pursuing projects centered around strong female characters. In fact, she emphasized that she intends to focus on creating films primarily featuring women rather than catering exclusively to male audiences.
\end{Verbatim}

\subsection{Example 3: test position
82}\label{record-docs-watermarking-qualitative-paper-sir-bira-md-example-3-test-position-82}

Prompt ID:
\texttt{\seqsplit{a8aa4f7effda3b9a9f5d3ee789d1d3c1dcb65a4df31ca926e955d1241abda5e1}}.

\textbf{Question / prompt}

\begin{Verbatim}
A 35-year-old cop with the Railway Protection Force had to be given four sutures on the head after a Nigerian drug peddler
\end{Verbatim}

\textbf{Original watermarked-model output}

\begin{Verbatim}
A 35-year-old cop with the Railway Protection Force had to be given four sutures on the head after a Nigerian drug peddler smashed his head against an electricity conductor after he objected to his mobile phone being taken from him.
The incident happened here on April 1, when Inspector Sukhwant Singh of RPF was checking the movement of Railway protection officers at Rajindra Cross in Delhi when Doremus Abdada, 32, of Odisha, tried to take his mobile phone. Inspector Sukhwait tried to take the mobile phone from him and started telling him to allow them to take him through security.
The accused, however, smashed one of the policeman’s forehead and claimed he could not hear. As he couldn’t distinguish between words and sounds, Inspector Sukhshit Singh asked him to give a physical description. Mr. Doremus, however, said he could not give one.
The Inspector, said Senior Executive Magistrate Dharmendra Nath Singh, said, he had a problem imagining that he himself had been injured. Constable Manoj, a resident of Malviya Nagar, was driving Inspector Sukhwit Singh’s vehicle.
The police said they were patrolling Rajindra Crossing when they received a PCR call
\end{Verbatim}

\textbf{After / alternative output}

\begin{Verbatim}
A railway inspector with the Railway Protection Force sustained serious head injuries on April 1st when a suspect attempting to steal his phone physically assaulted him. The incident occurred at Rajindra Cross in Delhi as Inspector Sukhwant Singh was conducting routine checks on railway personnel. The suspect, identified as 32-year-old Doremus Abdada from Odisha, attempted to seize the inspector's phone. 

Upon being confronted by the inspector, who demanded cooperation for further questioning, Doremus instead violently struck Inspector Sukhwant's head with considerable force. When questioned about his actions due to claims of hearing difficulties, the suspect was unable to provide a coherent response. 

Inspector Sukhwant subsequently requested a detailed physical description for identification purposes; however, this request was met with further confusion as Doremus insisted on his inability to provide such information. According to eyewitness accounts provided by Senior Executive Magistrate Dharmendra Nath Singh, it became clear that the suspect was experiencing difficulty accepting his own injuries.

The officer who suffered these injuries was being driven by his colleague, Constable Manoj, at the time. The incident occurred during routine patrols at Rajindra Crossing where the team responded to a Police Control Room (PCR) call.
\end{Verbatim}

\subsection{Example 4: test position
123}\label{record-docs-watermarking-qualitative-paper-sir-bira-md-example-4-test-position-123}

Prompt ID:
\texttt{\seqsplit{adaf979fdfb964904e387f40af11efc473f86dfd03a188d448f24221cea6cf05}}.

\textbf{Question / prompt}

\begin{Verbatim}
"Samantha Who?" has emerged as one of the breakout hits of this troubled TV season, and a major component of the ABC sitcom's
\end{Verbatim}

\textbf{Original watermarked-model output}

\begin{Verbatim}
"Samantha Who?" has emerged as one of the breakout hits of this troubled TV season, and a major component of the ABC sitcom's overall ensemble. We got to talk with showrunner Tracey Wigfield on set in Los Angeles and get an update on its future, including its surprising and successful run and whether or not this character has any kind of a future beyond season two and its finale in March 2012...

Now, the question arises... how will Samantha take over her father's house in season three, how will Sam (ahem) and Samantha interact with her new friends [U+2014] and not to mention, how will Samantha be able to keep up her relationship with her ex Chris (ahem, Chris B.) and her mom? As Wigfield explains, "Samantha takes over from father Bob (David Farr) in season three." And not just any old season, but three, not including the original [U+2014] or original and two later [U+2014] airing Season One and its infamous cliffhanger, even if Samantha does take over.

When asked to predict when Samantha takes over, Wigfield offers, "I don't want to be too specific." She also wants to let us know something [U+2014] in regard to her character Samantha, who will apparently be
\end{Verbatim}

\textbf{After / alternative output}

\begin{Verbatim}
The standout success of "Samantha Who?" has been instrumental in defining this challenging television season, with the ABC sitcom's diverse cast playing a pivotal role within its framework. In conversation with showrunner Tracey Wigfield, who kindly shared insights while filming at their Los Angeles location, we gained valuable information regarding the series' trajectory, encompassing both its notable success thus far and potential developments for key characters beyond the second season's conclusion scheduled for March 2012.

Several pressing questions arise regarding upcoming storylines: How might Samantha assume control of her late father's residence during season three? Furthermore, what dynamics can viewers expect between Samantha, now known as Sam, as she establishes friendships, while also maintaining relationships with significant figures such as her mother, former partner Chris, whose full name is revealed as Chris B., and navigating complex interactions within these new connections? According to Wigfield, "Samantha assumes ownership of the property left behind by her deceased father, Bob, played by David Farr, during season three, marking yet another pivotal chapter within the narrative."

When prompted for specific details regarding these upcoming developments, Wigfield opted for discretion, stating, "I prefer to avoid providing overly precise predictions about these events, allowing room for creative evolution as storylines unfold naturally within their respective context, while also choosing to maintain some secrecy surrounding key plot points, particularly those related to Samantha's evolving journey, which promises much excitement for fans eagerly awaiting further episodes."
\end{Verbatim}

\subsection{Example 5: test position
165}\label{record-docs-watermarking-qualitative-paper-sir-bira-md-example-5-test-position-165}

Prompt ID:
\texttt{\seqsplit{1b3ba6e9a2d857aa37a67122f763b84f6e273823246f7c6fcec9fced3d89d443}}.

\textbf{Question / prompt}

\begin{Verbatim}
LUFKIN, TX (KTRE) - Sunday was a day several years in the making at Lufkin’s First Baptist Church
\end{Verbatim}

\textbf{Original watermarked-model output}

\begin{Verbatim}
LUFKIN, TX (KTRE) - Sunday was a day several years in the making at Lufkin’s First Baptist Church. The men and women who serve as pastors and members of First Baptist held their official service online through Zoom.
\end{Verbatim}

\textbf{After / alternative output}

\begin{Verbatim}
In a long-awaited moment for the community of Lufkin's First Baptist Church, a significant milestone took place on the church's premises. The occasion, marking a culmination of efforts spanning multiple years, saw the church's leaders, comprising both male and female clergy, conduct their inaugural virtual worship services via the video conferencing platform, Zoom.
\end{Verbatim}

\subsection{Example 6: test position
206}\label{record-docs-watermarking-qualitative-paper-sir-bira-md-example-6-test-position-206}

Prompt ID:
\texttt{\seqsplit{03347445733f407fb1ccdcdfee0ac8a6870be235b90c64510ee96dfa7af4dfcf}}.

\textbf{Question / prompt}

\begin{Verbatim}
Pulitzer Prize-winning author Philip Roth resorts to public channels to get a Wikipedia entry changed...and succeeds.
After writing an open (
\end{Verbatim}

\textbf{Original watermarked-model output}

\begin{Verbatim}
Pulitzer Prize-winning author Philip Roth resorts to public channels to get a Wikipedia entry changed...and succeeds.
After writing an open (and highly successful) essay contest for Wikipedia contributors to name his new book, Philip G. Roth, he wants to trademark the word, and has enlisted the assistance of some of the world's sharpest lawyers to take a shot at gaining control of its use in print and in the electronic medium. But just when he thinks he's got his ducks in a row, his project has its first, and big, problem. So how can he win back support from not only Wikipedia, but also from his legions of fans, many of whom have rallied around the challenge he laid on them, and from his former publisher, Houghton Mifflin Harcourt, where he has been editor since he departed from his previous editor, Neal Gabler, in 2012, after Gabler left to work at Penguin Group. We talk with Philip. [U+2014] Adam Clark
\end{Verbatim}

\textbf{After / alternative output}

\begin{Verbatim}
Prolific Pulitzer Prize-winning author Philip Roth employs unconventional means to alter a Wikipedia entry and ultimately achieves success in doing so. Following the popularity of his open essay contest for Wikipedia contributors to suggest the title of his upcoming book, 'Philip G. Roth,' Roth proceeds to seek trademark protection for the title with the help of top-notch legal professionals who aim to secure control over the title's use both in printed materials and digital platforms. However, Roth soon encounters an initial and significant setback that threatens to undermine the momentum of his initiative and jeopardize the backing of key stakeholders including the online community at large, as well as the devoted fan base that had enthusiastically responded to the original challenge and his former publishing company, Houghton Mifflin Harcourt, which had been home to Roth following his departure from the editorial team led by Neal Gabler in 2012; the latter having transitioned to the Penguin Group at the time of their parting ways.
\end{Verbatim}

\subsection{Example 7: test position
247}\label{record-docs-watermarking-qualitative-paper-sir-bira-md-example-7-test-position-247}

Prompt ID:
\texttt{\seqsplit{4cbf4e262824ea813199d26415d873e7a935d3d4f9e628144b1679a9ef93d63a}}.

\textbf{Question / prompt}

\begin{Verbatim}
After the box leaked online earlier this week, Lego has officially taken the wraps off its Simpsons set.
The Lego Simpsons set is coming, officially
\end{Verbatim}

\textbf{Original watermarked-model output}

\begin{Verbatim}
After the box leaked online earlier this week, Lego has officially taken the wraps off its Simpsons set.
The Lego Simpsons set is coming, officially. Set photos have surfaced online and will be coming soon with more details coming on official release online at 9am EST. The Simpsons will also make its television debut on an episode of The Simpsons this month, and will be making more TV appearances in 2019. But in preparation for these new appearances, the Simpsons creator Matt Groening has hinted at some Simpsons characters making appearances in movies soon, too. Check out previous Simpsons appearances in movies here and here, and see how he says it could play into movies here.
It's not the Simpsons you're looking for, but it could be something else! The Simpsons: The Movie will be in theaters in 2019 and while details are thin at this point, perhaps we'll get to see some other characters from the past make an appearance, or have some familiar characters make an unexpected appearance. So perhaps we'll see Uncle Moe from Happy Days, who would be fun to have. Or maybe we'll just get to see an old friend and maybe see The Chief Marge again, though we don't think he has an age limit. He could be in his 90s, if his voice has remained intact,
\end{Verbatim}

\textbf{After / alternative output}

\begin{Verbatim}
Lego is now openly showcasing its highly anticipated Simpsons-themed set following a leak that emerged online a few days ago. Official images of the upcoming collection have been made available and further information is expected to be released shortly via their website at 9am EST. As part of a promotional push for its return to television later this month, a character-based episode of The Simpsons is scheduled to air; moreover, additional TV appearances are planned for next year. In anticipation of its expanded media presence, Matt Groening, the creative mind behind the series, has made a subtle hint that several beloved Simpsons characters may be featured in upcoming films as well. For a glimpse of its cinematic history, readers can refer to previous movie roles played by various Simpson's personalities.

Attention is being diverted to a different subject for now; however, fans may find interest in a separate development involving The Simpsons franchise – specifically "The Simpsons Movie," which is scheduled to premiere in theaters next year despite sparse information available about its storyline as of now. There is a possibility that audiences may catch a glimpse of lesser-known or returning characters within its narrative; alternatively, established figures may surprise viewers by reappearing or showing up unexpectedly. Although not a certainty, potential candidates to appear within "The Simpsons Movie" might include a cameo by a character like "Uncle Moe" from "Happy Days" or a reemergence of "The Chief" character; albeit speculative suggestions, they remain intriguing possibilities for viewers to look forward to – particularly since "The Chief" is reportedly still eligible to appear despite being potentially over ninety years old as long as their vocal capabilities remain unaltered
\end{Verbatim}

\subsection{Example 8: test position
288}\label{record-docs-watermarking-qualitative-paper-sir-bira-md-example-8-test-position-288}

Prompt ID:
\texttt{\seqsplit{9cdf452ec7392837b8e59d71537c6bd8a5a860e4379361d145f1b68dc8ef065b}}.

\textbf{Question / prompt}

\begin{Verbatim}
Lincoln Chafee is running for the Democratic nomination for president in 2016. He served one four-year term as Rhode Island's governor,
\end{Verbatim}

\textbf{Original watermarked-model output}

\begin{Verbatim}
Lincoln Chafee is running for the Democratic nomination for president in 2016. He served one four-year term as Rhode Island's governor, from 2007 to 2011. He was governor in 2009, when Providence, which is home to a large Chinese diaspora, went completely $1 billion in the red from 2008 through 2009. Chafee said his experience in dealing with municipal governments and state budgets in Rhode Island made him an ideal candidate to represent Rhode Island in D.C.

Toby Ann Curry, the host of NPR's The Place to Be, returned to her hometown to share her two new books, The Place She Ate and The Place She Sang, and to talk about writing and her work with TED. From The New Yorker magazine, here are quotes from three authors and authors-to-be who spoke from and around NPR's studios: Chafee, Curry and her mother, Nobel Peace Prize laureate Malala Yousafzai.

Author Sarah Duhatschek has been featured in The Boston Globe, The New York Times and The Boston Herald, among a dozen other national and international publications for her book, The Last Mile: An Essay On Sports, Family, And Happiness On The Last Mile That Fails And Fails And Gets
\end{Verbatim}

\textbf{After / alternative output}

\begin{Verbatim}
Rhode Island native Lincoln Chafee sought the Democratic presidential nomination in 2016 after serving as the Ocean State's governor from 2007 to 2011. Specifically, he was at the helm during a challenging fiscal period; as of 2009, Providence's budget deficit reached nearly $1 billion over the course of just one year (2008-2009). Reflecting on this time as governor, Chafee emphasized that his hands-on experience managing local governments and state finances would serve him well as a representative of Rhode Island in Washington D.C.

Meanwhile, author Toby Ann Stavreti-Curry (host of NPR's 'The Place To Be') revisited her hometown to discuss both of her recently released books - 'The Place She Ate' and 'The Place She Sang' - as well as shed light on her creative process and involvement with TED (Technology Entertainment Design). Attending this event at NPR's studios were notable guests including Lincoln Chafee; Toby's mother; Malala Yousafzai - renowned peace activist and recipient of the prestigious Nobel Peace Prize; along with up-and-coming author Sarah Duhatschek whose notable essay collection 'The Last Mile' garnered significant attention across various prominent publications such as 'The Boston Globe', 'The New York Times', 'The Boston Herald', as well as numerous others globally
\end{Verbatim}

\href{records/docs/watermarking/text.tex}{Text metrics} \textbar{}
\href{../docs/watermarking/README.md}{Report overview}
